\documentclass[11pt,letterpaper]{article}
\usepackage[margin=.85in]{geometry}
\usepackage[T1]{fontenc}
\usepackage[utf8]{inputenc}
\usepackage{lmodern,microtype,amsmath,amssymb,amsthm,hyperref}
\hypersetup{colorlinks=true,linkcolor=blue,citecolor=blue,urlcolor=blue,
pdftitle={An odd-period pedal-antipedal area invariant in elliptic billiards},
pdfauthor={Alec Kriebel}}
\newtheorem{theorem}{Theorem}
\newtheorem{lemma}[theorem]{Lemma}
\newtheorem{remark}[theorem]{Remark}
\DeclareMathOperator{\sn}{sn}
\DeclareMathOperator{\cn}{cn}
\DeclareMathOperator{\dn}{dn}
\newcommand{\area}{\mathcal A}
\newcommand{\ii}{\mathrm i}
\title{An odd-period pedal-antipedal area invariant\\in elliptic billiards}
\author{Alec Kriebel\\
\small Independent researcher\\
\small\href{https://orcid.org/0009-0001-9320-500X}{ORCID: 0009-0001-9320-500X}}
\date{13 October 2026}
\begin{document}
\maketitle
\begin{abstract}
We prove invariant $k_{403,a}$ observed by Reznik, Garcia and Koiller:
for every primitive odd-period elliptic billiard family with a strict
confocal elliptical caustic, the signed areas of the original orbit's
origin pedal and origin antipedal have constant product. All admitted
coprime star windings are included, and every real antipedal intersection
is finite. A complete pole analysis distinguishes coincident from
opposite complex endpoints. Odd-period arithmetic then places the
antipedal area and the credited pedal trace in complementary
one-dimensional elliptic-function spaces. A direct argument covers
circles; an exact four-periodic example rules out dropping the odd
restriction.
\end{abstract}

\section{The source invariant and conventions}
For an ordered closed polygon $V=(V_0,\ldots,V_{N-1})$, put
\[
 \area(V)=\frac12\sum_{j=0}^{N-1}\det(V_j,V_{j+1}),
 \qquad V_N=V_0.
\]
This is signed shoelace area, including for stars and self-crossing
derived polygons. Translate the billiard ellipse's center $O$ to zero.
For its original orbit $P$, let $q_j$ be the perpendicular foot from
$O$ to the line $P_jP_{j+1}$. Its origin-antipedal vertices are
\[
 R_j=L_j\cap L_{j+1},\qquad
 L_j=\{X:P_j\cdot X=P_j\cdot P_j\}.
\]
Thus $L_j$ passes through the original vertex and is perpendicular to
its radius. Write $A_O=\area(q)$ and $A_O^*=\area(R)$.

Reznik, Garcia and Koiller recorded the conjectural product
$A_OA_O^*$ for odd $N$ in their 2020 table; the same unproved row
$k_{403,a}$ appears in the published Table~5
\cite{source,old}. Both constructions concern the original orbit.
We give a complete proof of that observation in the strict
elliptical-caustic setting.

\begin{theorem}\label{thm:main}
Fix a noncircular ellipse with semiaxes $a>b>0$ and a strictly nested,
nondegenerate confocal elliptical caustic. In any billiard family of
least odd period $N\geq3$, the product $A_OA_O^*$ is phase independent.
Every admitted coprime star winding is included. All consecutive
antipedal intersections are finite.
\end{theorem}
No division by a derived signed area is involved. Reversal negates
both areas and preserves their product. Repeating an admissible cycle
$r$ times multiplies the product by $r^2$; the hypothesis concerns the
least period, rather than parity manufactured by a repeated list.
Hyperbolic or collapsed caustics are outside the theorem.

The parametrization and complex-analytic billiard methods are classical \cite{stachel,dlmf,ast}. Cyclic Jacobi
trace reductions and their Landen representations are established
tools \cite{cyclic,landen}; we do not claim these methods or the
observed invariant as new. The origin-pedal mechanism was already
proved in the earlier campaign's $k_{203,b}$ argument \cite{pedal}.
We restate it here so that the proof is self-contained. The additional
step is the construction-specific antipedal divisor analysis,
including its opposite-endpoint poles. Bounded primary-source
priority audits found no earlier covering proof of this exact
odd-origin product, without certifying historical firstness.

\section{Canonical coordinates and a trace lemma}
Let the caustic semiaxes be $\alpha>\beta>0$, put
$k'=\beta/\alpha$, $k^2=1-k'^2$, and let $K,K'$ be the complete
elliptic integrals. All Jacobi functions below have modulus $k$.
Set $p=\ii K'$. Choose forward orientation and coprime winding $\tau$
so that $0<\tau<N/2$. Stachel's canonical theorem
\cite[Theorem~4.3 and (4.9)]{stachel} gives
\begin{equation}\label{eq:param}
 v=\frac{2K\tau}{N},\quad \delta=2v,\quad
 a=\frac{\alpha\dn v}{\cn v},\quad b=\frac{\beta}{\cn v},
 \quad P(w)=(-a\sn w,b\cn w),\quad P_j=P(w+j\delta).
\end{equation}
Here $0<v<K$ and $N\delta=4K\tau$. Write
$s=\sn v$, $c=\cn v$, $d=\dn v$, all positive. Addition formulas give
\begin{equation}\label{eq:B}
 B(u):=\det(P(u-v),P(u+v))
 =\frac{2absc\,\dn u}{D(u)},\qquad
 D(u)=1-k^2s^2\sn^2u.
\end{equation}
For real $u$, $\dn u>0$ and $D(u)\geq1-k^2s^2>0$.
Therefore consecutive real radial vectors are independent and every
$R_j$ exists and is finite, also for stars.

The classical shifts used below are \cite[22.4 and 22.8]{dlmf}
\begin{align}
 &\sn(z+2K)=-\sn z,\quad \cn(z+2K)=-\cn z,\quad
 \dn(z+2K)=\dn z, \label{eq:realshift}\\
 &\sn(z+2p)=\sn z,\quad \cn(z+2p)=-\cn z,\quad
 \dn(z+2p)=-\dn z, \label{eq:imagshift}\\
 &\sn(p+z)=\frac1{k\sn z},\quad
 \cn(p+z)=\frac{-\ii\dn z}{k\sn z},\quad
 \dn(p+z)=\frac{-\ii\cn z}{\sn z}. \label{eq:pchart}
\end{align}
Complex dot products are bilinear, with no conjugation.

Put $\ell=2K/N$ and
\[
 S(w)=\sum_{j=0}^{N-1}\dn(w+j\delta),\qquad
 X=\mathbb C/(\ell\mathbb Z+4p\mathbb Z).
\]
\begin{lemma}\label{lem:trace}
$S$ is even, $\ell$-periodic and $2p$-antiperiodic. On $X$ it has
exactly two simple poles, at $p,3p$, and exactly two simple zeros, at
$p+\ell/2,3p+\ell/2$. Consequently $S(w)S(w+\ell/2)$ is constant.
Every meromorphic function on $X$ with only possible simple poles at
$p,3p$ and antiperiod $2p$ is a scalar multiple of $S$.
\end{lemma}
\begin{proof}
Since $N$ is odd and $\gcd(N,\tau)=1$, multiplication by
$\delta/\ell=2\tau$ permutes the real subgroup of order $N$ modulo
$2K$. Reindexing gives period $\ell$ and evenness; \eqref{eq:imagshift}
gives the antiperiod. The simple poles of $\dn$ are
$2rK+(2t+1)p$. At $w=p$ only one summand is singular modulo $2K$,
with residue $-\ii$ by \eqref{eq:pchart}; it cannot cancel.
The antiperiod gives the opposite residue at $3p$. These exhaust
the quotient pole classes.

Negation followed by translation by $\ell+2p$ fixes
$p+\ell/2$, while changing the function's sign. This forces a zero;
the same argument gives the other displayed zero. They are distinct
nonpoles. The trace is not identically zero since it is positive on
the real axis. Equality of total zero and pole orders on a compact
torus now makes these its complete simple zero divisor.
The product with the half-real-period translate has no poles and is
holomorphic on $X$, hence constant.

For the final assertion, match the residue at $p$ to that of $S$.
The antiperiod matches the second residue automatically. The
difference is holomorphic on $X$, hence constant, and its antiperiod
forces that constant to be zero.
\end{proof}

\section{The actual antipedal area and every possible pole}
For complex vectors define $H(P)=P\cdot P$ and
\begin{equation}\label{eq:F}
 F(P,Q)=
 \frac{H(P)H(Q)-(H(P)+H(Q))(P\cdot Q)/2}{\det(P,Q)}.
\end{equation}
The actual antipedal signed area has the rational expression
\begin{equation}\label{eq:U}
 U(w):=A_O^*(w)=\sum_jF(P_j,P_{j+1}).
\end{equation}
To verify it directly, let $J(x,y)=(-y,x)$. The intersections
preceding and following $P_i$ are
\[
 \begin{aligned}
 R_i&=P_i+\frac{H(P_{i+1})-P_i\cdot P_{i+1}}
                      {\det(P_i,P_{i+1})}JP_i,\\
 R_{i-1}&=P_i-\frac{H(P_{i-1})-P_{i-1}\cdot P_i}
                      {\det(P_{i-1},P_i)}JP_i.
\end{aligned}
\]
They satisfy their named line equations. Expanding
$\det(R_{i-1},R_i)/2$ and collecting the two contributions belonging
to each original edge gives \eqref{eq:F}--\eqref{eq:U}.
There is no convexity assumption or division by a complex norm.

We classify the complete divisor support of \eqref{eq:B} modulo
$2K,2p$; the imaginary congruence classifies support, and is an
antiperiod rather than a period of $B$.
The elliptic function $\sn^2u$ has one double pole in this cell.
The two distinct roots of $D$ are $p\pm v$, with
$D'(p\pm v)=\pm2cd/s\ne0$. At these points $\dn$ is finite and
nonzero, so $B$ has simple poles, and exactly one chord endpoint
is a pole vertex. At $p$, $\dn$ has a simple pole and $D$ a double
pole, so $B$ has a simple zero. The finite endpoints are opposite:
$P(p-v)=-P(p+v)$. At $r=K+p$, $\dn$ has a simple zero and
$D(r)=c^2\ne0$, so $B$ has a simple zero with coincident finite
endpoints, using $P(2r-z)=P(z)$. There are no other zeros or poles.

There are therefore exactly three cases for \eqref{eq:U}.
At coincident finite endpoints, its numerator has the identity
\begin{equation}\label{eq:polarization}
\begin{aligned}
 &H(P)H(Q)-\tfrac12(H(P)+H(Q))(P\cdot Q)\\
 &\qquad=\tfrac14\bigl((H(P)+H(Q))H(P-Q)-(H(P)-H(Q))^2\bigr).
\end{aligned}
\end{equation}
It is $O(z^2)$ when $P-Q=O(z)$, while $B$ has a simple zero.
The apparent singularity is removable, even if a finite complex
norm is zero. At opposite finite endpoints the numerator is merely
regular and the simple determinant zero gives at most a simple
pole. This pole must be retained.

At a pole vertex $P(p+z)$ the two neighbors are regular, since
$0<\delta<2K$. Each incident $F$ has numerator order at most three
and denominator a simple pole, hence pole order at most two.
Define $G(z)=F(P(p+z),P(p+\delta+z))$.
The identities $P(p-z)=-P(p+z)$,
$P(p-\delta+z)=-P(p+\delta-z)$,
$F(Q,P)=-F(P,Q)$ and $F(-P,-Q)=F(P,Q)$ give the sum of its
incident terms exactly as
\begin{equation}\label{eq:cancellation}
 G(z)-G(-z).
\end{equation}
The even negative Laurent coefficient cancels, leaving at most a
simple pole. Odd coprimality ensures only one pole vertex per
such phase modulo $2K$, so no edge has two pole endpoints.
For $N=3$, the edge between its two finite neighbors is the separate
opposite-endpoint term already bounded above, and adds at most a
simple pole.

Vertex-pole classes in $w$ are $p-j\delta$ and $3p-j\delta$.
Opposite-midpoint pole classes are $p-v-j\delta$ and
$3p-v-j\delta$. Since $v=\tau\ell$, all reduce to $p,3p$ on $X$.
This is the indispensable odd-period arithmetic. The symmetries of
$F$ and \eqref{eq:realshift}--\eqref{eq:imagshift} make $U$
$2K$-periodic, cyclically $\delta$-periodic and $2p$-antiperiodic,
hence $\ell$-periodic. We have proved that it has no other poles,
and those allowed have order at most one.
Lemma~\ref{lem:trace} consequently gives
\begin{equation}\label{eq:antitrace}
 U(w)=c_*S(w)
\end{equation}
for a real constant $c_*$, which may be zero.

\section{The credited origin-pedal factor and conclusion}
For clarity we reproduce the needed earlier pedal argument
\cite{pedal}. The actual chord in \eqref{eq:B} is the caustic tangent
\[
 n(u)\cdot X=1,\qquad n(u)=(-\sn u/\alpha,\cn u/\beta).
\]
Both endpoints satisfy it by addition formulas. Its origin foot is
\begin{equation}\label{eq:foot}
 Q(u)=\frac{n(u)}{n(u)\cdot n(u)}
   =\left(-\frac{\alpha k'^2\sn u}{\dn^2u},
                  \frac{\beta\cn u}{\dn^2u}\right).
\end{equation}
This is finite for real $u$. Define
\[
 T(u)=\frac12\sum_j\det(Q(u+j\delta),Q(u+(j+1)\delta)).
\]
The actual original-side pedal area is $A_O(w)=T(w+v)$.

At common Jacobi poles $Q$ extends holomorphically and vanishes.
Its only possible coordinate poles are the simple zeros of $\dn$,
at $r=K+p$ and translates, with order at most two.
The shifts and parity give $Q(2r-z)=Q(z)$, so it is even about $r$.
The two incident area terms combine as
\[
 \frac12\det\bigl(Q(r+z),
                   Q(r+\delta+z)-Q(r-\delta+z)\bigr).
\]
The neighbors are regular; by the even reflection their difference
is $O(z)$. The resulting pole is at most simple.
Moreover $Q(u+2K)=-Q(u)$ and
$Q(u+2p)=\operatorname{diag}(1,-1)Q(u)$.
Thus $T$ is $\ell$-periodic and $2p$-antiperiodic, with only possible
simple poles at $r,r+2p$ on $X$. Applying the shifted uniqueness
assertion of Lemma~\ref{lem:trace} yields
\begin{equation}\label{eq:pedaltrace}
 T(u)=c_pS(u+K)
\end{equation}
with real $c_p$.

Combining \eqref{eq:antitrace} and \eqref{eq:pedaltrace} gives
\[
 A_O(w)A_O^*(w)=c_pc_*S(w+v+K)S(w).
\]
Since $(v+K)/\ell=\tau+N/2$ is a half-integer,
this is the constant half-shift product in Lemma~\ref{lem:trace}.
This proves Theorem~\ref{thm:main}, without dividing by either area.

For a circular billiard of radius $R$, a primitive forward star has
step $\theta=2\pi\tau/N\in(0,\pi)$. The two derived regular stars
have radii $R\cos(\theta/2)$ and $R/\cos(\theta/2)$.
Their finite-intersection signed area product is
\[
 \frac{N^2R^4}{4}\sin^2\theta.
\]
This treats the circular extension directly, without using a
degenerate elliptic torus.

\section{Exact controls and verification}
The odd restriction is substantive. For $a^2=5,b^2=3$, the strict
four-periodic caustic has squared semiaxes $25/8,9/8$.
The axis diamond and the rectangle with vertices
$(\pm\sqrt{25/8},\pm\sqrt{9/8})$ are primitive reflected trajectories
in that same family. Tangency and reflection follow from
$\alpha^2/a^2+\beta^2/b^2=1$ and
$\alpha/a^2=\beta/b^2$ for the rectangle. Direct projections and
antipedal intersections give products $225/4$ and $289/4$.
Thus an unrestricted even-period extension would be false.

The verification archive contains the manuscript source, explicit
line-construction and symbolic checks, exact reflected three- and
four-periodic certificates, primitive-star arithmetic controls,
independent real and near-complex-pole diagnostics, source hashes,
bounded priority reports and reproducibility instructions.
The analytic argument establishes the all-period theorem; finite
checks corroborate its algebra and domain. High-precision outputs
are diagnostics, not interval certificates.

AI tools were used extensively in solving, verification, literature
research, coding, writing and internal adversarial review.
This is an unrefereed preprint, without external human peer review
or proof-assistant formalization. The original observation and
prior mathematical inputs remain explicitly credited.

\begin{thebibliography}{9}
\small\raggedright
\bibitem{source}
D.~Reznik, R.~Garcia and J.~Koiller,
Fifty new invariants of $N$-periodics in the elliptic billiard,
\emph{Arnold Math. J.} \textbf{7} (2021), 341--355.
\href{https://doi.org/10.1007/s40598-021-00174-y}{doi:10.1007/s40598-021-00174-y}.
\bibitem{old}
D.~Reznik, R.~Garcia and J.~Koiller,
Eighty new invariants in the elliptic billiard,
\href{https://arxiv.org/abs/2004.12497v11}{arXiv:2004.12497v11} (2020),
Table~5, p.~7.
\bibitem{stachel}
H.~Stachel, On the motion of billiards in ellipses,
\emph{Eur. J. Math.} \textbf{8} (2022), 1602--1622.
\href{https://doi.org/10.1007/s40879-021-00524-2}{doi:10.1007/s40879-021-00524-2}.
\bibitem{dlmf}
NIST Digital Library of Mathematical Functions,
\href{https://dlmf.nist.gov/22.4}{Section~22.4},
\href{https://dlmf.nist.gov/22.8}{Section~22.8} and
\href{https://dlmf.nist.gov/22.9}{Section~22.9}
(periods, shifts, addition and cyclic identities); accessed 13 October 2026.
\bibitem{ast}
A.~Akopyan, R.~Schwartz and S.~Tabachnikov,
Billiards in ellipses revisited,
\emph{Eur. J. Math.} \textbf{8} (2022), 1313--1327.
\href{https://doi.org/10.1007/s40879-020-00426-9}{doi:10.1007/s40879-020-00426-9}.
\bibitem{cyclic}
A.~Khare, A.~Lakshminarayan and U.~Sukhatme,
Cyclic identities involving Jacobi elliptic functions. II,
\href{https://arxiv.org/abs/math-ph/0207019v2}{arXiv:math-ph/0207019v2} (2003),
Section~2.1.1.
\bibitem{landen}
A.~Khare and U.~Sukhatme,
Generalized Landen transformation formulas for Jacobi elliptic functions,
\href{https://arxiv.org/abs/math-ph/0312074v1}{arXiv:math-ph/0312074v1} (2003),
equations~(13)--(15).
\bibitem{pedal}
A.~Kriebel, The origin-pedal area product in an elliptic billiard,
campaign proof for $k_{203,b}$, 1 October 2026;
\href{https://github.com/AlecKriebel/Math/blob/551a524e2db65fec90684485ef676d0b4179378d/unsolved_math_prioritization/attempts/5100012/PROOF.md}{immutable reviewed PR204 proof}.
\end{thebibliography}
\end{document}
